\documentclass[norsk]{beamer}
 
\usepackage[T1]{fontenc}
\usepackage{textcomp,pdfpages}
\usepackage{babel}

\usepackage{amsmath, amsfonts, epsfig, xspace}
\usepackage{pstricks,pst-node}
\usepackage{multimedia}
\usepackage{beamerthemesplit}
\usepackage[absolute, overlay]{textpos}
\usepackage{listings}

\usepackage{color}
 
\definecolor{dkgreen}{rgb}{0,0.6,0}
\definecolor{gray}{rgb}{0.5,0.5,0.5}
\definecolor{mauve}{rgb}{0.58,0,0.82}

\lstset{
  language=java,
  aboveskip=3mm,
  belowskip=3mm,
  showstringspaces=false,
  columns=flexible,
  basicstyle={\tiny\ttfamily},
  numberstyle=\tiny\color{gray},
  keywordstyle=\color{blue},
  commentstyle=\color{dkgreen},
  stringstyle=\color{mauve},
  breaklines=true,
  breakatwhitespace=true,
  tabsize=2
}
 


\usetheme{bekk}


\author{Aslak Johannessen \\ aslakjo@bekk.no  @aslakjo}

\title[Scala]{Actors}
\subtitle[]{Scala -actors}
\institute{BEKK Consulting AS}


\begin{document}

\maketitle

\begin{frame}{Scala quiz}
  \begin{content}
    Oppgave: Lag et quizlag og implementer det som en actor\\
    
    \includegraphics[scale=0.9]{quizlag}
  \end{content}
\end{frame}

\begin{frame}{Scala quiz}
  \begin{content}
    Quiz p� vanlig m�te, med en liten vri\\
    \center
    1. Be om et sp�rsm�l og f�r sp�rsm�l\\
    \includegraphics[scale=0.6]{quiz-sporsmalsvar}\\
    2. Svar p� sp�rsm�let og f�r vite om det er riktig
  \end{content}
\end{frame}


\begin{frame}{Actors - hva er greia}
  \begin{content}
    Wikipedia: \\
    \begin{quotation}In computer science, the Actor model is a mathematical model of concurrent computation that treats "actors" as the universal primitives of concurrent digital computation: in response to a message that it receives, an actor can make local decisions, create more actors, send more messages, and determine how to respond to the next message received. The Actor model originated in 1973. It has been used both as a framework for a theoretical understanding of concurrency, and as the theoretical basis for several practical implementations of concurrent systems.
    \end{quotation}
  \end{content}
\end{frame}

\begin{frame}{Actors - hva er greia}
  \begin{content}
    Wikipedia: \\
    \begin{quotation}... \emph{concurrent computation} ... \emph{in response to a message that it receives, an actor can make local decisions ... and        determine how to respond to the next message received.}  \emph{The Actor model originated in 1973.} ...
    \end{quotation}
  \end{content}
\end{frame}


\begin{frame}[fragile]{Actors - eksempler eksempler}
  \begin{content}
    \begin{lstlisting}[basicstyle={\large\ttfamily}]
      case class Tick
      
      class Counter extends Actor {
        private var counter = 0

        def receive = {
          case Tick =>
            counter += 1
            println(counter)
        }
      }
    \end{lstlisting}
  \end{content}
\end{frame}

\begin{frame}{Actors - in action}
  \begin{content}
    \centering
    1. MoreQuestions(new Team(``lag en''))\\
    2. Question(``Ping'', Nil)\\
    \begin{columns}
      \column{.1\textwidth}
        Server
      \column{.5\textwidth}
        \includegraphics[scale=0.5]{quiz-sporsmalsvar}
      \column{.1\textwidth}
        Lag
    \end{columns}
    3. Answer(question, ``pong'')\\
    4. Correct()
  \end{content}
\end{frame}

\begin{frame}[containsverbatim]{Actors - in action}
  \begin{content}
    \centering
    \begin{lstlisting}[basicstyle={\small\ttfamily}]
case class MoreQuestions(val team:Team)
case class Question(val question:String, val content:Any)
    \end{lstlisting}
    \begin{columns}
      \column{.1\textwidth}
        Server
      \column{.5\textwidth}
        \includegraphics[scale=0.5]{quiz-sporsmalsvar}
      \column{.1\textwidth}
        Lag
    \end{columns}
    \begin{lstlisting}[basicstyle={\small\ttfamily}]
case class Answer(val team:Team, val question:Question, val answer:Any)
case class Correct extends Verdict
    \end{lstlisting}
  \end{content}
\end{frame}

\begin{frame}[containsverbatim]{Actors - in action}
  \begin{content}
    \begin{columns}
      \column{.35\textwidth}
        \includegraphics[scale=0.3]{quiz-sporsmalsvar}
      \column{.8\textwidth}
        \begin{lstlisting}[basicstyle={\small\ttfamily}]
val question = 
  (remote !! MoreQuestions(team))
  .get.asInstanceOf[Question]
remote !! new Answer( team, question, giveAnswer(question) )

def giveAnswer(quiestion: Question)={
  question match {
    case x@Question("Ping", _) =>  "pong"
  }
}
        \end{lstlisting}
    \end{columns}
  \end{content}
\end{frame}


\begin{frame}[containsverbatim]{Oppgave - meldinger}
  \begin{content}
    \begin{lstlisting}
      package no.bekk.scala.messages
      trait Message
      trait Verdict

      case class MoreChallenges(val team:Team) extends Message
      case class Question(val question:String, val content:Any) extends Message

      case class Answer(val temaName:Team, val question:Question, val answer: Any) extends Message
      
      case class Correct() extends Message with Verdict 
      case class Wrong() extends Message with Verdict
    \end{lstlisting}
  \end{content}
\end{frame}


\begin{frame}[containsverbatim]{Oppgave - meldinger}
  \begin{content}
    \begin{lstlisting}
      package no.bekk.scala.messages
      trait Message
      trait Verdict

      case class MoreChallenges(val team:Team) extends Message
      case class Question(val question:String, val content:Any) extends Message

      case class Answer(val temaName:Team, val question:Question, val answer: Any) extends Message
      
      case class Correct() extends Message with Verdict 
      case class Wrong() extends Message with Verdict
    \end{lstlisting}
  \end{content}

  �nsker �;
    \begin{itemize}
      \item ta i mot et sp�rsm�l Question, 
      \item finne ut hva det handler om, og 
      \item besvare
    \end{itemize}
\end{frame}

\kontrast{Oppgave \\\vfill\small https://github.com/aslakjo/quiz-client}


\kontrast{Actors}

\begin{frame}{Actors - hva er det}
  \huge
  \begin{content}
    \begin{equation*}
      \left. 
        \begin{aligned}
          \text{naturlig programerings model}\\
          \text{enkelt � paralellisere}\\
          \text{gode biblioteker}\\
          \text{enkelt � gj�re riktig} 
        \end{aligned} 
      \right\}
        Actors
    \end{equation*}
  \end{content}
\end{frame}

\kontrast{Hvordan?}

\begin{frame}{Actors - hva er det}
  \huge
  \begin{content}
    \begin{equation*}
      \left.
        Actors
      \right.
        \begin{cases}
          &\text{individer er normalt}\\
          &\text{ingen deling}\\
          &\text{akka actors erlang}\\
          &\text{objekter} 
        \end{cases} 
    \end{equation*}
  \end{content}
\end{frame}

\begin{frame}{Actors - helt enkelt}
  \begin{content}
    \includegraphics[width=\textwidth]{mailbox}\\
    Ingen deling \hfill-\hfill meldingsboks \hfill-\hfill eksterne
  \end{content}
\end{frame}

\begin{frame}[containsverbatim]{Actors - Akka}
  \begin{content}
    \begin{lstlisting}
      import se.scalablesolutions.akka.actor._
      import se.scalablesolutions.akka.remote._

      class Client extends GiveAnswer
      {
        val team = new Team("test")
        val remote = RemoteClient.actorFor("Server", "localhost", 9999)
        
        def run {

          while(true)
          {
            val question = (remote !! MoreChallenges(team)).get.asInstanceOf[Question]
            remote !! new Answer(team, question, giveAnswer(question) )
          }
        }

      }
    \end{lstlisting}
  \end{content}
\end{frame}

\kontrast{Videre - oppgave}



\kontrast{Pattern matching}

\begin{frame}[containsverbatim]{Pattern matching - enkelt}
  \begin{content}
    \begin{lstlisting}[basicstyle={\large\ttfamily}]
      newMessage match {
	      case Tick => counter += 1
	      case Tack(i:Int) => counter += i
	      case message@Clear => counter = 0
	      case "add a hundred" => counter += 100
	      case 2 =>counter += 2
	      case _ => // Do nothing
      }
    \end{lstlisting} 
  \end{content}
\end{frame}

\begin{frame}[containsverbatim]{Pattern matching - enkelt}
  \centering
  \begin{content}
    Type pattern
    \begin{lstlisting}[basicstyle={\large\ttfamily}]
	      case Tick => counter += 1
    \end{lstlisting} 
    Matcher p� typen, Tick
  \end{content}
\end{frame}

\begin{frame}[containsverbatim]{Pattern matching - enkelt}
  \centering
  \begin{content}
    Constructor pattern
    \begin{lstlisting}[basicstyle={\large\ttfamily}]
	      case Tack(i:Int) => counter += i
    \end{lstlisting} 
    Matcher p� typen, Tick, med hvor den interne variablen er en Int
  \end{content}
\end{frame}

\begin{frame}[containsverbatim]{Pattern matching - enkelt}
  \centering
  \begin{content}
    Variable binding pattern
    \begin{lstlisting}[basicstyle={\large\ttfamily}]
	      case message@Clear => counter = 0
        case message:Clear => counter = 0 //compile err.
    \end{lstlisting} 
    Henter ut Clear objektet til message variablen.
  \end{content}
\end{frame}

\begin{frame}[containsverbatim]{Pattern matching - enkelt}
  \centering
  \begin{content}
    Constant pattern
    \begin{lstlisting}[basicstyle={\large\ttfamily}]
	      case  "add a hundred" => counter += 100
	      case 2 =>counter += 2
    \end{lstlisting} 
    Matcher singletons som er lik seg selv.
  \end{content}
\end{frame}

\begin{frame}[containsverbatim]{Pattern matching - enkelt}
  \centering
  \begin{content}
    Wildcard pattern
    \begin{lstlisting}[basicstyle={\large\ttfamily}]
	      case _ => // Do nothing
    \end{lstlisting} 
    Matcher alt.
  \end{content}
\end{frame}

\kontrast{Oppgaver}

\begin{frame}[containsverbatim]{Pattern matching - avansert}
  \begin{content}
    Matche p� sp�rsm�l\\
    Pakke opp objecter, extractors\\
    \begin{lstlisting}[basicstyle={\ttfamily}]
new Question("sp�rsm�l", List("s", "v")) match {
  case x@Question(sp, svar:List[String]) =>  {
    println(sporsmal + ", " + svar)//warning!
  }
}
    \end{lstlisting}
  \end{content}
\end{frame}

\begin{frame}[containsverbatim]{Pattern matching - avansert}
  \begin{content}
    Matche p� sp�rsm�l\\
    Pakke opp objecter, extractors\\
    \begin{lstlisting}[basicstyle={\ttfamily}]
new Question("sp�rsm�l", List("s", "v")) match {
  case x@Question(sp, svar@List(s:String, _*) =>{
   println(sporsmal + ", " + svar)
  }
}
    \end{lstlisting}
  \end{content}
\end{frame}

\begin{frame}[containsverbatim]{Pattern matching - avansert}
  \begin{content}
    Matche p� sp�rsm�l\\
    Regexp\\
    \begin{lstlisting}[basicstyle={\large\ttfamily}]
	    val Epost = "^(.*)@(.*)\.(.*)$".r
    \end{lstlisting}
  \end{content}
\end{frame}

\begin{frame}[containsverbatim]{Pattern matching - avansert}
  \begin{content}
    Matche p� sp�rsm�l\\
    Regexp\\
    \begin{lstlisting}[basicstyle={\large\ttfamily}]
	    val Epost = "^(.*)@(.*)\.(.*)$".r

      "aslakjo@gmail.com" match {
	      case Epost(user, domain, tld) => //spam!
      }
    \end{lstlisting}
  \end{content}
\end{frame}

\begin{frame}[plain]
  \begin{centering}
    \includegraphics[height=11em]{bekk}
    \\\vspace{2em}\hfill\\
    \small
    \includegraphics[height=1em]{bekk-logo}
    \\\vspace{.4em}\hfill\\
    \fontsize{8}{8}
    \selectfont
    \textbf{Aslak Johannessen}\\
    Senior Consultant\\
    982 19 249 \\
    aslakjo@bekk.no\\
    \vspace{.3em}\hfill\\
    \fontsize{1}{1}
    \selectfont
    BEKK CONSULTING AS\\
    SKUR 39, VIPPETANGEN. P.O. BOX 134 SENTRUM, 0102 OSLO, NORWAY. WWW.BEKK.NO\\
  \end{centering}
\end{frame}

\end{document}
